\chapter{The Slopes of $U_3$}
\label{chap:jacobs}

We now instantiate the general theory at the case worked out by Jacobs \cite{Jacobs}, and prove
the theorem the thesis is aimed at: for the definite quaternion algebra over $\Q$ ramified at
$2$ and $\infty$, at level $U_1(9)$ and at a suitable weight, the eigenvalues of $U_3$ have
$3$-adic valuations
\[
  \tfrac12, \quad \tfrac32, \quad \tfrac52, \quad \dots
\]
--- an arithmetic progression of half-integers.  The half-integrality is not an artefact: it is
what makes the example interesting, and it is possible only because the coefficient field is
ramified over $\Q_3$.

Everything general has been done.  What remains is to exhibit the data --- the algebra, the
place, the level, the weight, the coset representatives --- and to compute one matrix.

\section{The data}

\begin{definition}
  \label{def:jacobs-setting}
  \lean{JacobsSlash.D, JacobsSlash.v₃, JacobsSlash.K₃, JacobsSlash.theta}
  \leanok
  Let $D = \bbH_{\Q}$ be Hamilton's quaternions over $\Q$, a definite quaternion algebra of
  discriminant $2$; let $v_3$ be the place above $3$ and $K_3 = \Q_3$ its completion.  Since $3
  \nmid 2$, $D$ splits at $v_3$, and a splitting is fixed by choosing a square root $\nu$ of
  $-2$ in $K_3$; the resulting isomorphism $\theta_3 : D \otimes_{\Q} K_3 \cong M_2(K_3)$ is the
  rigidification of Definition~\ref{def:rigidification}.
\end{definition}

%% ν₃ is located explicitly by Hensel at 2695 mod 3^10; the location is what makes the matrix of
%% U₃ computable rather than merely existent.

\begin{definition}
  \label{def:hurwitz}
  \lean{JacobsSlash.hurwitzOrder}
  \leanok
  The \emph{Hurwitz order} is $\calO = \Z\langle i, j, \tfrac12(1+i+j+k)\rangle \subseteq \bbH_{\Q}$,
  the unique maximal order up to conjugacy.  We record membership by the parity condition that
  $2x$ has integer coordinates all of the same parity.
\end{definition}

\begin{proposition}
  \label{prop:hurwitz-units}
  \uses{def:hurwitz}
  \lean{JacobsSlash.isUnit_iff_hnorm_eq_one, JacobsSlash.card_units_hurwitzOrder}
  \leanok
  A Hurwitz quaternion is a unit exactly when its reduced norm is $1$, and $\calO^\times$ has
  exactly $24$ elements.
\end{proposition}

\begin{definition}
  \label{def:jacobs-levels}
  \uses{def:hurwitz, def:sigma0}
  \lean{JacobsSlash.U0, JacobsSlash.U1_9}
  \leanok
  $U_0(1) \subseteq D_f^\times$ consists of the adelic units integral at every finite place ---
  the unit group of $\calO \otimes \Zhat$ --- and $U_1(9)$ refines it by the congruence
  $\theta_3(g) \equiv \left(\begin{smallmatrix} * & * \\ 0 & 1\end{smallmatrix}\right)$ modulo
  $9$.
\end{definition}

\begin{proposition}
  \label{prop:jacobs-levels-compact-open}
  \uses{def:jacobs-levels, prop:compact-open-levels}
  \lean{JacobsSlash.isCompact_U0, JacobsSlash.isCompact_U1_9, JacobsSlash.isOpen_U1_9,
        JacobsSlash.finite_image_doubleCoset_U1_9}
  \leanok
  $U_0(1)$ and $U_1(9)$ are compact open, so Theorem~\ref{thm:hecke-finiteness} applies and the
  Hecke operators at this level are defined without exhibiting representatives.
\end{proposition}

\begin{proof}
  \leanok
  \uses{prop:compact-open-levels}
  The everywhere-integral elements of $D \otimes \A_f$ are exactly the $\Zhat$-span of the
  Hurwitz basis $1, i, j, \omega$: one inclusion is clear, and the other extracts each coordinate
  by a $\Q$-linear functional, extends it to the completion, and uses that a rational integral at
  every finite place is a rational integer.  That set is compact by
  Proposition~\ref{prop:compact-open-levels} and open because the integral adeles are open, so
  $U_0(1)$ is compact open; $U_1(9)$ is cut out of it by a congruence condition at $3$, which is
  clopen because closed balls in an ultrametric field are open.
\end{proof}

\section{Class number one}

\begin{theorem}[{[Jacobs, Theorem 2.1]}]
  \label{thm:jacobs-classset}
  \uses{def:jacobs-levels, thm:fujisaki}
  \lean{JacobsSlash.hClassNumberOne, JacobsSlash.classRep_bijective}
  \leanok
  $D_f^\times = D^\times \cdot U_0(1)$, and the class set at level $U_1(9)$ has exactly three
  elements, represented by explicit adelic $c_0, c_1, c_2$.
\end{theorem}

\begin{proof}
  \leanok
  Class number one for the Hurwitz order is classical; formalised here by the idelic dictionary,
  reducing a global integrality statement to local approximation.  Passing from $U_0(1)$ to
  $U_1(9)$ replaces the class set by the orbits of $\calO^\times$ on the primitive vectors of
  $(\Z/9)^2$, which are computed to be three.
\end{proof}

\begin{lemma}[{[Jacobs, Lemma 2.2]}]
  \label{lem:jacobs-neat}
  \uses{thm:jacobs-classset, thm:evalatreps}
  \lean{JacobsSlash.stabilizerAt_classRep}
  \leanok
  The stabilisers $\Gamma_i = c_i^{-1}D^\times c_i \cap U_1(9)$ are trivial.
\end{lemma}

So the level is neat, and Theorem~\ref{thm:evalatreps} gives $S_\kappa(U_1(9)) \cong c(\mathrm{Fin}\,3
\times \N, K_3)$: a form is exactly a triple of elements of the Tate algebra.

\section{The weight}

\begin{definition}
  \label{def:jacobs-weight}
  \uses{def:analyticweight}
  \lean{JacobsSlash.jacobsCharOf, JacobsSlash.jacobsWeight}
  \leanok
  For $t \in K_3$ with $\lVert t \rVert < 1$ let $\kappa_t(u) = u^t := \exp(t \log u)$ on the
  $1$-units.  The expansion of $\kappa_t(cz+d)$ is the binomial series, whose coefficients
  satisfy the decay required by Definition~\ref{def:weightseries}; the pair is an analytic weight
  in the sense of Definition~\ref{def:analyticweight}.
\end{definition}

%% The hypothesis ‖t‖ < 1 is Jacobs' t ∈ m₃ (p. 38), not an artefact of the formalisation.

\section{The operator and its matrix}

\begin{definition}
  \label{def:u3}
  \uses{def:jacobs-weight, def:heckeoperator, def:jacobs-levels}
  \lean{JacobsSlash.eta3, JacobsSlash.heckeU3}
  \leanok
  $U_3 = [U_1(9)\, \eta_3\, U_1(9)]$ where $\eta_3$ is the adelic element which is $1$ away from
  $3$ and $\left(\begin{smallmatrix}1&0\\0&3\end{smallmatrix}\right)$ at $3$.
\end{definition}

\begin{theorem}
  \label{thm:u3-matrix}
  \uses{def:u3, thm:transport, lem:jacobs-neat}
  \lean{JacobsSlash.heckeU3_apply_classRep, JacobsSlash.matrixCoeff_blockEntry}
  \leanok
  With respect to the three representatives, $(U_3\varphi)(c_i) = \sum_j \varepsilon_{ij}
  (\varphi(c_j))$, and the matrix coefficients of the blocks $\varepsilon_{ij}$ are the explicit
  coefficients of the generating functions of Proposition~\ref{prop:kappaslash-matrix} evaluated
  at the certificate elements.
\end{theorem}

The certificate data --- four right coset representatives $v_t$ of $U_1(9)\eta_3 U_1(9)$, the
permutation $\sigma$, and the factorisations $c_i v_t^{-1} = d\, c_{\sigma(i,t)}\, u$ --- is
computed by hand; it is the only part of the argument that is a computation rather than a
theorem, and it is what the general framework was built to consume.

\begin{definition}
  \label{def:u3-det}
  \uses{thm:u3-matrix, def:heckecharpowerseries}
  \lean{JacobsSlash.charPowerSeriesU3}
  \leanok
  $\det(1 - T\,U_3)$, computed in the block model through the three evaluations.
\end{definition}

\section{The eigenvalues}

Over $K_3$ the matrix of $U_3$ does not have unit top-left minors, and no arithmetic progression
of integer slopes can appear: the answer involves half-integers, so the polygon must be read over
a ramified extension.  Adjoining a primitive cube root of unity $\omega$ gives $K_3(\omega) =
\Q_3(\sqrt{-3})$, ramified of degree $2$, in which $v_3$ takes half-integer values; there the
matrix splits into eigenblocks.

\begin{theorem}
  \label{thm:u3-factorisation}
  \uses{def:u3-det, prop:qmf-basechange}
  \lean{JacobsSlash.charPowerSeriesU3_factorisation}
  \leanok
  Over an isometric extension containing $\omega$, $\det(1 - TU_3)$ factors according to the
  $\omega$-eigenspaces of the block structure, and on the relevant factor the rescaled matrix has
  every top-left minor of norm $1$.
\end{theorem}

\begin{theorem}[The crux]
  \label{thm:u3-eigenvalues}
  \uses{thm:u3-factorisation, thm:unit-minors, cor:slope-reading}
  \lean{JacobsSlash.exists_eigenvalue_U3_halfIntegral}
  \leanok
  For every $j \in \N$ there is an eigenvalue $a$ of $U_3$ with
  \[
    v_3(a) = j + \tfrac12 .
  \]
\end{theorem}

\begin{proof}
  \leanok
  \uses{thm:unit-minors, cor:slope-reading, thm:np-slope-reading}
  Theorem~\ref{thm:u3-factorisation} puts the relevant factor in the shape of
  Theorem~\ref{thm:unit-minors} with $\varpi = \sqrt{-3}$, whose valuation is $\tfrac12$ in the
  normalisation $v_3(3) = 1$.  The exact coefficient valuations $\binom m2$ in that normalisation
  become $\tfrac12\binom m2$ here, whose successive differences are $\tfrac12, \tfrac32, \dots$;
  Corollary~\ref{cor:slope-reading} reads these off the polygon and
  Theorem~\ref{thm:serre12} converts slopes into eigenvalues.
\end{proof}

\begin{theorem}[The crux, for forms]
  \label{thm:u3-eigenforms}
  \uses{thm:u3-eigenvalues, thm:qmf-eigenform-criterion, prop:qmf-basechange}
  \lean{JacobsSlash.exists_eigenform_U3_halfIntegral}
  \leanok
  For every $j \in \N$ there exist $a \in \bbC_3$ and a nonzero overconvergent quaternionic
  modular form $\varphi$ of weight $\kappa_t$ and level $U_1(9)$ over $\bbC_3$ with
  \[
    U_3 \varphi = a\varphi, \qquad v_3(a) = j + \tfrac12 .
  \]
\end{theorem}

\begin{proof}
  \leanok
  \uses{thm:qmf-eigenform-criterion, thm:evalatreps}
  Theorem~\ref{thm:u3-eigenvalues} produces an eigenvector in the block model; the level is neat
  by Lemma~\ref{lem:jacobs-neat}, so Theorem~\ref{thm:evalatreps} transports it back to a form,
  and the transport intertwines the block operator with $U_3$ by Theorem~\ref{thm:transport}.
  Base change along $K_3(\omega) \hookrightarrow \bbC_3$ is
  Proposition~\ref{prop:qmf-basechange}.
\end{proof}

This is the statement announced in the introduction: the slopes of a compact Hecke operator, for
this weight and level, form an arithmetic progression --- and one that is visibly not integral,
so that the fibre of the ``quaternionic eigenvariety'' over this weight is not what a naive guess
would predict.
